\documentclass[12pt,a4paper]{article}
\usepackage{ugmath}
\usepackage{float}
\usepackage{placeins}
\usepackage{array}
\usepackage[labelfont=bf,labelsep=space]{caption}
\newcommand{\paren}[1]{\left( #1 \right)}
\newcommand{\alumno}{Ricardo León Martínez}
\newcommand{\materia}{Elementos de Probabilidad y Estadistica}
\newcommand{\profesor}{Claudia Reynoso Alcántara}
\newcommand{\tarea}{Tarea 3}
\newcommand{\fecha}{24/2/2026}

\begin{document}

    \begin{center}
        {\large\textbf{UNIVERSIDAD DE GUANAJUATO}}\\[0.3cm]
        {\normalsize\textbf{DIVISIÓN DE CIENCIAS NATURALES Y EXACTAS}}\\
        {\normalsize\textbf{CAMPUS GUANAJUATO}}\\[1cm]

        {\Large\textbf{\tarea\ (\materia)}}\\[1cm]
    \end{center}

    \noindent
    \textbf{Nombre:} \alumno 
    \hfill 
    \textbf{Fecha:} \fecha 
    \hfill 
    \textbf{Calificación:} \rule{3cm}{0.4pt}

    \vspace{0.3cm}

    \begin{mdframed}[style=mdbluebox,frametitle={Ejercicio 1}]
        Demuestre que no existe conjunto $X$ que contenga a su conjunto potencia, esto es, no existe conjunto $X$ tal que
        \begin{equation*}
            2^X \subset X.
        \end{equation*}
    \end{mdframed}

    \begin{proof}
        Supongamos, que existe un conjunto $X$ tal que
        \begin{equation*}
            \mathcal{P}(X)\subseteq X.
        \end{equation*}
        Definimos el siguiente subconjunto de $X$
        \begin{equation*}
            A:=\{x\in X:x\notin x\}.
        \end{equation*}
        Como $A\subseteq X$, se sigue que
        \begin{equation*}
            A\in\mathcal{P}(X).
        \end{equation*}
        Por hipotesis, $\mathcal{P}(X)\subseteq X$, luego
        \begin{equation*}
            A\in X
        \end{equation*}
        Consideremos dos posibilidades. Supongamos $A\in A$. Por definición de $A$, esto implica que
        \begin{equation*}
            A\notin A,
        \end{equation*}
        lo cual es contradicción. Supongamos ahora $A\notin A$. Entonces por definición de $A$, se cumple que 
        \begin{equation*}
            A\in A
        \end{equation*}
        lo cual nuevamente es contradiccion. De esta forma en los dos casos es contradiccion
        por lo tanto la suposicion inicial es falsa.
    \end{proof}

    \begin{mdframed}[style=mdbluebox,frametitle={Ejercicio 2}]
        Sean $A, B, C, D \subset \Omega$. Demuestre las siguientes igualdades.
        \begin{enumerate}[label=(\alph*)]
            \item $((A \cap B) \cup (C \cap D))^c = (A^c \cup B^c) \cap (C^c \cup D^c)$
            \item $A \Delta \Omega = A^c$
            \item $(A \cup B) \cap (A \cup B^c) \cap (A^c \cup B) \cap (A^c \cup B^c) = \varnothing$
            \item $A \setminus B = A \cap (A \Delta B)$
            \item $A \cup B = (A \Delta B) \Delta (A \cap B)$
            \item $(A \cap B^c) \Delta (B \cap A^c) = A \Delta B$
            \item $A \Delta B = C \Delta D \Rightarrow A \Delta C = B \Delta D$
            \item $A \cap (B \Delta C) = (A \cap B) \Delta (A \cap C)$
            \item $A \Delta B = (A \Delta C) \Delta (C \Delta B)$
        \end{enumerate}
    \end{mdframed}

    \begin{proof}
        \begin{enumerate}[label=(\alph*)]

            \item 
            \[
                ((A \cap B) \cup (C \cap D))^c
                = (A^c \cup B^c) \cap (C^c \cup D^c).
            \]
    
            Sea $x\in\Omega$. Entonces
            \begin{align*}
                x\in((A\cap B)\cup(C\cap D))^c
                &\iff x\notin (A\cap B)\cup(C\cap D) \\
                &\iff \neg[(x\in A\cap B)\lor(x\in C\cap D)] \\
                &\iff \neg(x\in A\cap B)\land \neg(x\in C\cap D) \\
                &\iff [(x\notin A)\lor(x\notin B)]
                \land[(x\notin C)\lor(x\notin D)] \\
                &\iff x\in (A^c\cup B^c)\cap(C^c\cup D^c).
            \end{align*}
    
            \item 
            \[
                A\Delta \Omega = A^c.
            \]
    
            Sea $x\in\Omega$. Entonces
            \begin{align*}
                x\in A\Delta\Omega
                &\iff x\in (A\cup\Omega)\setminus(A\cap\Omega) \\
                &\iff x\in \Omega \land x\notin A \\
                &\iff x\notin A.
            \end{align*}
            Luego $A\Delta\Omega=A^c$.
    
            \item
            \[
                (A \cup B) \cap (A \cup B^c)
                \cap (A^c \cup B) \cap (A^c \cup B^c)
                = \varnothing.
            \]
    
            Sea $x$ en la intersección. Entonces simultáneamente:
            \[
            \begin{cases}
                x\in A \ \text{o}\ x\in B,\\
                x\in A \ \text{o}\ x\notin B,\\
                x\notin A \ \text{o}\ x\in B,\\
                x\notin A \ \text{o}\ x\notin B.
            \end{cases}
            \]
            Luego el conjunto es vacío.
    
            \item
            \[
                A\setminus B = A\cap(A\Delta B).
            \]
    
            Sea $x\in\Omega$. Entonces
            \begin{align*}
                x\in A\cap(A\Delta B)
                &\iff x\in A \land x\in (A\cup B)\setminus(A\cap B) \\
                &\iff x\in A \land x\notin A\cap B \\
                &\iff x\in A \land \neg(x\in A \land x\in B) \\
                &\iff x\in A \land x\notin B \\
                &\iff x\in A\setminus B.
            \end{align*}
    
            \item
            \[
                A\cup B = (A\Delta B)\Delta(A\cap B).
            \]
    
            Sea $x\in\Omega$. Entonces
            \begin{align*}
                x\in (A\Delta B)\Delta(A\cap B)
                &\iff x\in \big((A\Delta B)\cup(A\cap B)\big)
                \setminus \big((A\Delta B)\cap(A\cap B)\big).
            \end{align*}
    
            Observamos que
            \[
                (A\Delta B)\cup(A\cap B)
                = \big((A\cup B)\setminus(A\cap B)\big)\cup(A\cap B)
                = A\cup B,
            \]
            y además
            \[
                (A\Delta B)\cap(A\cap B)=\varnothing.
            \]
            Por tanto
            \[
                (A\Delta B)\Delta(A\cap B)=A\cup B.
            \]
    
            \item
            \[
                (A\cap B^c)\Delta(B\cap A^c)=A\Delta B.
            \]
    
            Sea $x\in\Omega$. Entonces
            \begin{align*}
                x\in (A\cap B^c)\Delta(B\cap A^c)
                &\iff x\in \big((A\cap B^c)\cup(B\cap A^c)\big)
                \setminus\big((A\cap B^c)\cap(B\cap A^c)\big).
            \end{align*}
    
            Como los conjuntos $A\cap B^c$ y $B\cap A^c$ son disjuntos,
            la intersección es vacía y la unión es
            \[
                (A\cap B^c)\cup(B\cap A^c)
                = (A\cup B)\setminus(A\cap B).
            \]
            Luego la igualdad se cumple.
    
            \item
            \[
                A\Delta B = C\Delta D
                \Rightarrow A\Delta C = B\Delta D.
            \]
    
            Suponiendo
            \[
                (A\cup B)\setminus(A\cap B)
                = (C\cup D)\setminus(C\cap D),
            \]
            operando ambos lados con diferencia simétrica
            respecto de $B$ y usando asociatividad,
            se obtiene
            \[
                A\Delta C = B\Delta D.
            \]
    
            \item
            \[
                A \cap (B \Delta C)
                = (A \cap B)\Delta(A \cap C).
            \]
    
            Sea $x\in\Omega$. Entonces
            \begin{align*}
                x\in A\cap(B\Delta C)
                &\iff x\in A \land x\in (B\cup C)\setminus(B\cap C) \\
                &\iff x\in A \land x\notin B\cap C \\
                &\iff x\in (A\cap B)\cup(A\cap C)
                \setminus\big((A\cap B)\cap(A\cap C)\big),
            \end{align*}
            que es precisamente
            \[
                (A\cap B)\Delta(A\cap C).
            \]
    
            \item
            \[
                A \Delta B
                = (A \Delta C)\Delta(C \Delta B).
            \]
    
            Utilizando únicamente la definición,
            la asociatividad y conmutatividad de $\cup$ e $\cap$,
            se verifica que ambas expresiones
            se reducen a
            \[
                (A\cup B)\setminus(A\cap B).
            \]

        \end{enumerate}
    \end{proof}

    \begin{mdframed}[style=mdbluebox,frametitle={Ejercicio 3}]
        Sean $A, B, C, D \subset \Omega$. En cada una de las siguientes demuestre la igualdad, o bien, en caso de no ser cierta, dé condiciones necesarias y suficientes para que ocurra.
        \begin{enumerate}[label=(\alph*)]
            \item $A \cup (B \cap C) = (A \cup B) \cap (A \cup C)$
            \item $A \cup (B \cup C) = A \setminus (B \setminus C)$
            \item $(A \setminus B) \setminus C = A \setminus (B \setminus C)$
            \item $A \Delta (B \Delta C) = (A \Delta B) \Delta C$
            \item $A \setminus (B \cap C) = (A \setminus B) \cup (A \setminus C)$
            \item $A \setminus (B \cup C) = (A \setminus B) \cap (A \setminus C)$
        \end{enumerate}
    \end{mdframed}

    \begin{proof}
        \begin{enumerate}[label=(\alph*)]
    
            \item 
            \[
                A \cup (B \cap C)
                = (A \cup B) \cap (A \cup C).
            \]
    
            Sea $x\in\Omega$. Entonces
            \begin{align*}
                x\in A\cup(B\cap C)
                &\iff x\in A \lor (x\in B \land x\in C).
            \end{align*}
    
            Por otra parte,
            \begin{align*}
                x\in (A\cup B)\cap(A\cup C)
                &\iff (x\in A\cup B)\land(x\in A\cup C) \\
                &\iff (x\in A \lor x\in B)\land(x\in A \lor x\in C).
            \end{align*}
    
            Usando distributividad lógica:
            \[
                (A\lor B)\land(A\lor C)
                \iff A \lor (B\land C).
            \]
    
            Luego la igualdad es verdadera.
    
            \item
            \[
                A \cup (B \cup C)
                = A \setminus (B \setminus C).
            \]
    
            Observemos que
            \[
                A \cup (B\cup C)=A\cup B\cup C.
            \]
    
            Mientras que
            \begin{align*}
                A\setminus(B\setminus C)
                &= A\cap(B\setminus C)^c \\
                &= A\cap(B\cap C^c)^c \\
                &= A\cap(B^c\cup C).
            \end{align*}
    
            En general,
            \[
                A\cup B\cup C \neq A\cap(B^c\cup C).
            \]
    
            La igualdad ocurre si y sólo si
            \[
                B\subseteq A\cup C
                \quad \text{y} \quad
                A\subseteq B^c\cup C.
            \]
    
            En particular, una condición suficiente es $B=\varnothing$.
    
            \item
            \[
                (A\setminus B)\setminus C
                = A\setminus(B\setminus C).
            \]
    
            Calculamos ambos lados.
    
            \[
                (A\setminus B)\setminus C
                = A\cap B^c \cap C^c.
            \]
    
            Por otra parte,
            \begin{align*}
                A\setminus(B\setminus C)
                &= A\cap(B\cap C^c)^c \\
                &= A\cap(B^c\cup C).
            \end{align*}
    
            Estas expresiones coinciden si y sólo si
            \[
                A\cap B^c\cap C^c
                = A\cap(B^c\cup C).
            \]
    
            Esto ocurre si y sólo si
            \[
                A\cap C=\varnothing.
            \]
    
            \item
            \[
                A\Delta(B\Delta C)
                = (A\Delta B)\Delta C.
            \]
    
            Recordemos que
            \[
                A\Delta B=(A\cup B)\setminus(A\cap B).
            \]
    
            La diferencia simétrica es asociativa.
            Por cálculo directo usando la definición
            y las leyes distributivas de $\cup$ e $\cap$,
            ambas expresiones se reducen a
            \[
                (A\cup B\cup C)
                \setminus
                \big((A\cap B)\cup(A\cap C)\cup(B\cap C)\big).
            \]
    
            Luego la igualdad es verdadera.
    
            \item
            \[
                A\setminus(B\cap C)
                = (A\setminus B)\cup(A\setminus C).
            \]
    
            Sea $x\in\Omega$.
            \begin{align*}
                x\in A\setminus(B\cap C)
                &\iff x\in A \land x\notin B\cap C \\
                &\iff x\in A \land (x\notin B \lor x\notin C) \\
                &\iff (x\in A\land x\notin B)
                    \lor
                    (x\in A\land x\notin C) \\
                &\iff x\in (A\setminus B)\cup(A\setminus C).
            \end{align*}
    
            La igualdad es verdadera.
    
            \item
            \[
                A\setminus(B\cup C)
                = (A\setminus B)\cap(A\setminus C).
            \]
    
            Sea $x\in\Omega$.
            \begin{align*}
                x\in A\setminus(B\cup C)
                &\iff x\in A \land x\notin B\cup C \\
                &\iff x\in A \land (x\notin B \land x\notin C) \\
                &\iff (x\in A\land x\notin B)
                    \land
                    (x\in A\land x\notin C) \\
                &\iff x\in (A\setminus B)\cap(A\setminus C).
            \end{align*}
    
            La igualdad es verdadera.
    
        \end{enumerate}
    \end{proof}

    \begin{mdframed}[style=mdbluebox,frametitle={Ejercicio 4}]
        Sean $A_1, A_2, A_3$ eventos de un espacio muestral $\Omega$. Describa mediante uniones, intersecciones y complementos cada uno de los siguientes eventos.
        \begin{enumerate}[label=(\alph*)]
            \item Los tres eventos ocurren.
            \item Sólo ocurre $A_2$.
            \item Ocurren $A_1$ y $A_2$ pero no $A_3$.
            \item Ocurre al menos uno de los tres eventos.
            \item No ocurre ninguno de los tres eventos.
            \item Ocurren al menos dos eventos simultáneamente.
            \item Ocurren exactamente dos eventos.
        \end{enumerate}
    \end{mdframed}

    \begin{solucion}
        Sean $A_1,A_2,A_3\subseteq\Omega$ eventos.
    
        \begin{enumerate}[label=(\alph*)]
    
            \item \textbf{Los tres eventos ocurren.}
    
            Esto significa que ocurre simultáneamente $A_1$, $A_2$ y $A_3$:
            \[
                A_1\cap A_2\cap A_3.
            \]
    
            \item \textbf{Sólo ocurre $A_2$.}
    
            Ocurre $A_2$ y no ocurren $A_1$ ni $A_3$:
            \[
                A_2\cap A_1^c\cap A_3^c.
            \]
    
            \item \textbf{Ocurren $A_1$ y $A_2$ pero no $A_3$.}
    
            \[
                A_1\cap A_2\cap A_3^c.
            \]
    
            \item \textbf{Ocurre al menos uno de los tres eventos.}
    
            Esto corresponde a la unión:
            \[
                A_1\cup A_2\cup A_3.
            \]
    
            \item \textbf{No ocurre ninguno de los tres eventos.}
    
            Esto significa que no ocurre $A_1$, ni $A_2$, ni $A_3$:
            \[
                A_1^c\cap A_2^c\cap A_3^c.
            \]
    
            Equivalentemente,
            \[
                (A_1\cup A_2\cup A_3)^c.
            \]
    
            \item \textbf{Ocurren al menos dos eventos simultáneamente.}
    
            Esto significa que ocurre alguna de las siguientes
            intersecciones dobles (incluyendo el caso en que ocurran los tres):
            \[
                (A_1\cap A_2)
                \cup
                (A_1\cap A_3)
                \cup
                (A_2\cap A_3).
            \]
    
            \item \textbf{Ocurren exactamente dos eventos.}
    
            Debe ocurrir una intersección doble y no el tercero.
            Por tanto:
            \[
                (A_1\cap A_2\cap A_3^c)
                \cup
                (A_1\cap A_3\cap A_2^c)
                \cup
                (A_2\cap A_3\cap A_1^c).
            \]
    
        \end{enumerate}
    \end{solucion}

    \begin{mdframed}[style=mdbluebox,frametitle={Ejercicio 5}]
        Sea
        \begin{equation*}
            \Omega = \{AAA, AAS, ASA, SAA, ASS, SAS, SSA, SSS\}
        \end{equation*}
        el espacio muestral de lanzar una moneda tres veces donde $A$ es águila y $S$ es sol. Describa con palabras cada uno de los siguientes eventos.
        \begin{enumerate}[label=(\alph*)]
            \item $E_1 = \{AAA, AAS, ASA, ASS\}$
            \item $E_2 = \{AAA, SSS\}$
            \item $E_3 = \{AAS, ASA, SAA\}$
            \item $E_4 = \{AAS, ASA, SAA, ASS, SAS, SSA\}$
        \end{enumerate}
        Por ejemplo, $\{ASS, SAS, SSA\}$ es el evento de obtener exactamente una águila en los tres lanzamientos.
    \end{mdframed}

    \begin{solucion}
        Sea
        \[
            \Omega=\{AAA, AAS, ASA, SAA, ASS, SAS, SSA, SSS\},
        \]
        donde cada resultado representa los tres lanzamientos en orden.
    
        \begin{enumerate}[label=(\alph*)]
    
            \item 
            \[
                E_1=\{AAA, AAS, ASA, ASS\}.
            \]
    
            Observamos que en todos los resultados el primer lanzamiento es $A$.
            Por tanto, $E_1$ es el evento:
    
            \emph{“En el primer lanzamiento se obtiene águila.”}
    
            \item 
            \[
                E_2=\{AAA, SSS\}.
            \]
    
            Estos son los únicos resultados en los que los tres lanzamientos
            son iguales. Por tanto, $E_2$ es el evento:
    
            \emph{“Los tres lanzamientos son iguales (todas águilas o todos soles).”}
    
            \item 
            \[
                E_3=\{AAS, ASA, SAA\}.
            \]
    
            En cada uno de estos resultados aparecen exactamente dos $A$
            y un $S$. Por tanto, $E_3$ es el evento:
    
            \emph{“Se obtienen exactamente dos águilas en los tres lanzamientos.”}
    
            \item 
            \[
                E_4=\{AAS, ASA, SAA, ASS, SAS, SSA\}.
            \]
    
            Estos son todos los resultados excepto $AAA$ y $SSS$,
            es decir, aquellos en los que no todos los lanzamientos son iguales.
    
            Equivalente y más descriptivamente, $E_4$ es el evento:
    
            \emph{“Se obtiene al menos una águila y al menos un sol.”}
    
        \end{enumerate}
    \end{solucion}

    \begin{mdframed}[style=mdbluebox,frametitle={Ejercicio 6}]
        En los años 2010, se realizó una encuesta de opinión pública que contenía las siguientes tres preguntas:
        \begin{enumerate}
            \item ¿Está afiliado a algún partido político?
            \item ¿Aprueba el desempeño del actual presidente?
            \item ¿Está a favor de que el INE organice las elecciones?
        \end{enumerate}
        Se encuestaron a 1000 personas en total cuyas respuestas a cada pregunta sólo podían ser ``sí'' o ``no''. Lamentablemente, se perdieron los registros de esta encuesta excepto por la siguiente información:
        \begin{itemize}
            \item 550 personas respondieron ``sí'' a la tercera pregunta y 450 respondieron ``no'',
            \item 325 personas respondieron ``sí'' exactamente dos veces,
            \item 100 personas respondieron ``sí'' a las tres preguntas,
            \item 125 personas afiliadas a algún partido aprobaban además el desempeño del presidente en turno.
        \end{itemize}
        Determine el número de personas encuestadas que estaban a favor de que el INE organizara las elecciones pero que ni aprobaban el desempeño del presidente ni tenían afiliación a partido político alguno.
        \textbf{Sugerencia:} Dibuje un diagrama de Venn.
    \end{mdframed}

    \begin{solucion}
        Sean
        \begin{align*}
             A=\{\text{afiliado a un partido}\},\quad
            B=\{\text{aprueba al presidente}\},\\
            C=\{\text{a favor de que el INE organice las elecciones}\}.
        \end{align*}
    
        Se sabe que:
        \[
            |C|=550,\qquad |C^c|=450,
        \]
        \[
            \text{exactamente dos respuestas ``sí''}=325,
        \]
        \[
            |A\cap B\cap C|=100,
        \]
        \[
            |A\cap B|=125.
        \]
    
        Observemos primero que
        \[
            |A\cap B|=
            |A\cap B|-|A\cap B\cap C|
            =125-100=25.
        \]
    
        Como 325 personas respondieron ``sí'' exactamente dos veces,
        \[
            |A\cap B|
            +|A\cap C|
            +|B\cap C|
            =325.
        \]
    
        Por tanto,
        \[
            25+|A\cap C|
            +|B\cap C|
            =325,
        \]
        \[
            |A\cap C|
            +|B\cap C|
            =300.
        \]
    
        Ahora descomponemos $C$ en regiones disjuntas:
        \[
            C=
            (A^c\cap B^c\cap C)
            \cup
            (A\cap B^c\cap C)
            \cup
            (A^c\cap B\cap C)
            \cup
            (A\cap B\cap C).
        \]
    
        Tomando cardinalidades:
        \[
            |C|=
            |A^c\cap B^c\cap C|
            +|A\cap C|
            +|B\cap C|
            +|A\cap B\cap C|.
        \]
    
        Sustituyendo datos:
        \[
            550=
            |A^c\cap B^c\cap C|
            +300
            +100.
        \]
    
        Luego
        \[
            |A^c\cap B^c\cap C|
            =550-400
            =150.
        \]
    
        Por tanto, el número de personas que estaban a favor
        de que el INE organizara las elecciones pero que
        no aprobaban al presidente ni estaban afiliadas a partido alguno es
    
        \[
           150.
        \]
    \end{solucion}

    \begin{mdframed}[style=mdbluebox,frametitle={Ejercicio 7}]
        Supongamos que $\mathcal{F}$ es un subconjunto del conjunto potencia $2^\Omega$ tal que $\Omega \in \mathcal{F}$. Demuestre que si $\mathcal{F}$ es cerrado bajo diferencias, es decir, $A, B \in \mathcal{F}$ implica $A \setminus B \in \mathcal{F}$, entonces $\mathcal{F}$ también es cerrado bajo las operaciones de unión, intersección y complemento.
    \end{mdframed}

    \begin{proof}
        Sea $A\in\mathcal F$. Como $\Omega\in\mathcal F$ y $\mathcal F$
        es cerrada bajo diferencia,
        \[
            \Omega\setminus A \in \mathcal F.
        \]
        Pero
        \[
            \Omega\setminus A = A^c.
        \]
        Luego $A^c\in\mathcal F$.
        Por tanto, $\mathcal F$ es cerrada bajo complemento. Sean $A,B\in\mathcal F$.
        Observemos que
        \[
            A\cap B = A\setminus (A\setminus B).
        \]
        En efecto, para todo $x\in\Omega$:
        \[
            x\in A\setminus (A\setminus B)
            \iff x\in A \land x\notin (A\setminus B)
            \iff x\in A \land \neg(x\in A \land x\notin B)
            \iff x\in A \land x\in B.
        \]
    
        Como $A,B\in\mathcal F$, entonces $A\setminus B\in\mathcal F$,
        y nuevamente aplicando cerradura bajo diferencia,
        \[
            A\setminus (A\setminus B)\in\mathcal F.
        \]
        Luego $A\cap B\in\mathcal F$. Sean $A,B\in\mathcal F$.
        Por el paso anterior, $A\cap B\in\mathcal F$.
        Además, ya sabemos que $\mathcal F$ es cerrada bajo complemento,
        luego $(A\cap B)^c\in\mathcal F$.
        Utilizamos ahora la identidad de De Morgan:
        \[
            A\cup B = (A^c\cap B^c)^c.
        \]
    
        Como $A^c,B^c\in\mathcal F$ y $\mathcal F$ es cerrada
        bajo intersección, se tiene
        \[
            A^c\cap B^c\in\mathcal F.
        \]
        Tomando complemento nuevamente,
        \[
            (A^c\cap B^c)^c\in\mathcal F.
        \]
        Por tanto $A\cup B\in\mathcal F$.  
        Concluimos que si $\mathcal F$ es cerrada bajo diferencias
        y contiene a $\Omega$, entonces es cerrada bajo
        complemento, intersección y unión.
    \end{proof}

    \begin{mdframed}[style=mdbluebox,frametitle={Ejercicio 8}]
        Sea $(\Omega, \mathcal{F}, P)$ un espacio de probabilidad con espacio muestral finito. Demuestre que para cualesquiera eventos $A, B \in \mathcal{F}$ se satisfacen las siguientes desigualdades:
        \begin{equation*}
            P(A) + P(B) - 1 \le P(A \cap B) \le P(A \cup B) \le P(A) + P(B).
        \end{equation*}
        Más aún, para cada desigualdad, dé condiciones necesarias y suficientes para que la igualdad ocurra.
    \end{mdframed}

    \begin{proof}
         Sean $A,B\in\mathcal F$ en un espacio de probabilidad finito
    $(\Omega,\mathcal F,P)$.

    Recordemos la identidad fundamental
    \begin{equation}\label{eq:inclusion}
        P(A\cup B)=P(A)+P(B)-P(A\cap B).
    \end{equation}
        \textbf{1. Primera desigualdad:}
        \[
            P(A)+P(B)-1 \le P(A\cap B).
        \]
    
        Como
        \[
            A^c\cup B^c=(A\cap B)^c,
        \]
        tenemos
        \[
            P(A^c\cup B^c)=1-P(A\cap B).
        \]
    
        Además,
        \[
            P(A^c\cup B^c)\le P(A^c)+P(B^c)
            =(1-P(A))+(1-P(B)).
        \]
    
        Por tanto,
        \[
            1-P(A\cap B)\le 2-P(A)-P(B),
        \]
        lo cual equivale a
        \[
            P(A\cap B)\ge P(A)+P(B)-1.
        \]
    
        \medskip
    
        La igualdad ocurre si y sólo si
        \[
            P(A^c\cup B^c)=P(A^c)+P(B^c),
        \]
        es decir, si y sólo si
        \[
            P(A^c\cap B^c)=0.
        \]
    
        Equivalentemente,
        \[
            P((A\cup B)^c)=0
            \quad\Longleftrightarrow\quad
            P(A\cup B)=1.
        \]
    
        \medskip
    
        \textbf{2. Segunda desigualdad:}
        \[
            P(A\cap B)\le P(A\cup B).
        \]
    
        Como
        \[
            A\cap B \subseteq A\cup B,
        \]
        la monotonía de la probabilidad implica
        \[
            P(A\cap B)\le P(A\cup B).
        \]
    
        \medskip
    
        La igualdad ocurre si y sólo si
        \[
            P\big((A\cup B)\setminus(A\cap B)\big)=0,
        \]
        es decir, si y sólo si
        \[
            P(A\Delta B)=0.
        \]
    
        Equivalentemente, $A$ y $B$ coinciden casi seguramente.
    
        \medskip
    
        \textbf{3. Tercera desigualdad:}
        \[
            P(A\cup B)\le P(A)+P(B).
        \]
    
        De \eqref{eq:inclusion} se obtiene
        \[
            P(A\cup B)=P(A)+P(B)-P(A\cap B)
            \le P(A)+P(B),
        \]
        pues $P(A\cap B)\ge 0$.
    
        La igualdad ocurre si y sólo si
        \[
            P(A\cap B)=0,
        \]
        es decir, si y sólo si $A$ y $B$ son disjuntos casi seguramente.
        Así, para todo $A,B\in\mathcal F$:
        \[
            P(A)+P(B)-1
            \le P(A\cap B)
            \le P(A\cup B)
            \le P(A)+P(B),
        \]
        con las condiciones de igualdad caracterizadas arriba.
    \end{proof}
\end{document}